% =====================================================================
\chapter{Tổng quan về an toàn và bảo mật thông tin}
\label{chap:tongquan}
% =====================================================================

\section{Các khái niệm cơ bản}

Mật mã học (\tn{cryptography}) là ngành khoa học nghiên cứu các kỹ thuật bảo vệ thông tin
bằng cách biến đổi dữ liệu từ dạng có thể đọc được sang dạng không thể hiểu được đối với
người không có thẩm quyền. Mục đích của mật mã học là bảo đảm tính bí mật, tính toàn vẹn
và tính xác thực của thông tin trong quá trình lưu trữ cũng như truyền tải
\cite{stallings2020, menezes1996}. Ngày nay, mật mã học hiện diện trong hầu hết hoạt động
số: từ giao thức HTTPS khi duyệt web, giao dịch ngân hàng trực tuyến cho đến việc chia sẻ
hình ảnh trên mạng xã hội.

\subsection{Mã hóa và giải mã}

\textbf{Mã hóa} (\tn{encryption}) là quá trình dùng một thuật toán $E$ và một khóa $K$ để
biến bản rõ (\tn{plaintext}) $P$ thành bản mã (\tn{ciphertext}) $C$. \textbf{Giải mã}
(\tn{decryption}) là quá trình ngược lại, dùng thuật toán $D$ để khôi phục bản rõ từ bản mã:
\begin{equation}
  C = E_K(P), \qquad P = D_{K'}(C).
  \label{eq:enc-dec}
\end{equation}
Một hệ mật hoạt động đúng phải thỏa điều kiện $D_{K'}\bigl(E_K(P)\bigr) = P$ với mọi bản rõ
$P$. Nếu $K' = K$ ta có hệ mật khóa đối xứng; nếu $K' \neq K$ và không thể suy ra $K'$ từ
$K$ một cách hiệu quả, ta có hệ mật khóa bất đối xứng (khóa công khai).

\subsection{Khóa mã hóa và nguyên lý Kerckhoffs}

Khóa là thông tin bí mật quyết định phép biến đổi cụ thể của thuật toán. Theo
\textbf{nguyên lý Kerckhoffs} \cite{kerckhoffs1883}, độ an toàn của một hệ mật chỉ được phép
dựa vào bí mật của khóa, không được dựa vào việc giữ kín thuật toán. Hệ quả trực tiếp là
không gian khóa phải đủ lớn để chống tấn công vét cạn và khóa phải được sinh, lưu trữ, phân
phối an toàn. Nguyên lý này sẽ được dùng ở Chương~\ref{chap:vandung} để đánh giá lớp mã hóa
Chaotic của bài báo, vốn có các tham số được công bố công khai.

\section{Các mục tiêu của an toàn thông tin}

An toàn thông tin hướng đến bảo đảm ba thuộc tính cốt lõi thường gọi là bộ ba \textbf{CIA},
cùng với thuộc tính chống chối bỏ \cite{stallings2020}:
\begin{itemize}
  \item \textbf{Tính bí mật} (\tn{Confidentiality}): thông tin chỉ được truy cập bởi những
        người có thẩm quyền. Mã hóa là công cụ chính để đạt được tính bí mật.
  \item \textbf{Tính toàn vẹn} (\tn{Integrity}): thông tin không bị sửa đổi trái phép, hoặc
        nếu bị sửa đổi thì phải phát hiện được. Công cụ thường dùng là hàm băm mật mã, mã xác
        thực thông điệp (MAC) và chữ ký số.
  \item \textbf{Tính sẵn sàng} (\tn{Availability}): hệ thống và dữ liệu luôn sẵn sàng phục
        vụ người dùng hợp pháp khi cần; mối đe dọa điển hình là tấn công từ chối dịch vụ.
  \item \textbf{Tính chống chối bỏ} (\tn{Non-repudiation}): người gửi không thể phủ nhận đã
        gửi và người nhận không thể phủ nhận đã nhận thông tin; công cụ chính là chữ ký số.
\end{itemize}

Bên cạnh bộ ba CIA, trong bài toán truyền ảnh còn có một yêu cầu đặc thù là \textbf{tính
không thể nhận biết} (\tn{undetectability}): đối phương không được biết rằng có một thông
điệp bí mật đang được truyền đi. Đây chính là mục tiêu của kỹ thuật giấu tin trình bày ở
mục~\ref{sec:giautin}.

\section{Các loại thuật toán mật mã}

\subsection{Mã hóa đối xứng}

Mã hóa đối xứng dùng cùng một khóa bí mật cho cả mã hóa và giải mã. Hai bên liên lạc phải
chia sẻ trước khóa này qua một kênh an toàn. Ưu điểm lớn nhất của mã hóa đối xứng là tốc độ
cao, phù hợp để mã hóa khối lượng dữ liệu lớn như hình ảnh hay video \cite{stallings2020}.
Các thuật toán đối xứng chia thành hai nhóm:
\begin{itemize}
  \item \textbf{Mã khối} (\tn{block cipher}): xử lý dữ liệu theo từng khối có kích thước cố
        định, ví dụ DES và 3DES (khối 64 bit), AES (khối 128 bit). Để mã hóa thông điệp dài
        hơn một khối, mã khối được dùng kèm một \textbf{chế độ hoạt động} (\tn{mode of
        operation}) như ECB, CBC, CTR hay GCM \cite{sp80038a}.
  \item \textbf{Mã dòng} (\tn{stream cipher}): sinh ra một dòng khóa (\tn{keystream}) giả
        ngẫu nhiên rồi XOR với bản rõ theo từng byte hoặc từng bit, ví dụ RC4, ChaCha20. Các
        lược đồ mã hóa ảnh dựa trên hệ hỗn loạn cũng thường có cấu trúc của mã dòng.
\end{itemize}

Chế độ hoạt động ảnh hưởng trực tiếp đến độ an toàn khi mã hóa ảnh. Với chế độ ECB, các khối
bản rõ giống nhau cho ra các khối bản mã giống nhau nên đường nét của ảnh vẫn lộ ra trên bản
mã. Chế độ CBC khắc phục điều này bằng cách XOR mỗi khối bản rõ với khối bản mã liền trước
trước khi mã hóa; đây cũng là chế độ mà bài báo được tái hiện lựa chọn cho cả 3DES và AES.

\subsection{Mã hóa bất đối xứng}

Mã hóa bất đối xứng dùng một cặp khóa: khóa công khai để mã hóa (hoặc xác minh chữ ký) và
khóa riêng để giải mã (hoặc ký). Hệ RSA dựa trên độ khó của bài toán phân tích số nguyên lớn
ra thừa số nguyên tố, còn \textbf{mật mã đường cong elliptic} (ECC) dựa trên bài toán logarit
rời rạc trên nhóm điểm của đường cong elliptic. Với cùng mức an toàn, ECC cần khóa ngắn hơn
RSA nhiều lần \cite{katz2020}. Mã hóa bất đối xứng chậm hơn đáng kể so với mã đối xứng nên
trong thực tế thường chỉ dùng để trao đổi khóa phiên và ký số. Bài báo gốc có nhắc đến ECC,
nhưng như sẽ phân tích ở Chương~\ref{chap:baibao}, thành phần ``ECC'' trong bài báo thực chất
chỉ là một mặt nạ giả ngẫu nhiên được \emph{mô phỏng}, không có phép toán đường cong nào.

\subsection{Hàm băm mật mã}

Hàm băm mật mã $H$ ánh xạ đầu vào có độ dài tùy ý thành giá trị băm có độ dài cố định
(ví dụ 256 bit với SHA-256). Một hàm băm mật mã tốt có các tính chất:
\begin{itemize}
  \item \textbf{Một chiều}: biết $h$, không thể tìm được $m$ sao cho $H(m) = h$;
  \item \textbf{Kháng va chạm}: rất khó tìm hai đầu vào $m_1 \neq m_2$ với $H(m_1) = H(m_2)$;
  \item \textbf{Hiệu ứng tuyết lở}: thay đổi một bit đầu vào làm khoảng một nửa số bit đầu ra
        thay đổi.
\end{itemize}
Trong đề tài này, SHA-256 được dùng để \emph{dẫn xuất khóa} một cách tất định từ hạt giống
cấu hình hoặc từ mật khẩu người dùng (mục~\ref{sec:caidat}). Cần phân biệt hàm băm mật mã với
mã kiểm tra dư vòng CRC: CRC chỉ phát hiện lỗi ngẫu nhiên do kênh truyền gây ra, \emph{không}
chống được sửa đổi có chủ đích vì kẻ tấn công có thể tính lại CRC.

\section{Giấu tin trong ảnh số}
\label{sec:giautin}

\subsection{Khái niệm và mô hình}

Giấu tin (\tn{steganography}) là kỹ thuật nhúng một thông điệp bí mật vào một vật mang
(\tn{cover}) như ảnh, âm thanh hay video sao cho sự tồn tại của thông điệp không bị phát
hiện \cite{johnson1998, cheddad2010}. Mô hình kinh điển của giấu tin là ``bài toán người tù''
của Simmons \cite{simmons1984}: hai tù nhân muốn trao đổi kế hoạch vượt ngục qua người cai
ngục; mọi thư từ đều bị đọc, và chỉ cần một dấu hiệu khả nghi là liên lạc bị cắt. Vì vậy
thông điệp phải được giấu trong những lá thư trông hoàn toàn bình thường. Các thành phần của
một hệ giấu tin được minh họa ở \hinh{fig:stego-model}.

\begin{figure}[H]
\centering
\resizebox{\linewidth}{!}{%
\begin{tikzpicture}[node distance=4mm and 6mm, every node/.append style={font=\footnotesize}]
  \node[box] (msg) {Thông điệp\\bí mật};
  \node[box, above=of msg] (cover) {Ảnh phủ\\(\tn{cover})};
  \node[box, below=of msg] (key) {Khóa giấu tin\\(\tn{stego-key})};
  \node[pbox, right=9mm of msg, minimum height=13mm] (emb) {Bộ nhúng\\(\tn{embedding})};
  \node[box, right=of emb] (stego) {Ảnh mang tin\\(\tn{stego})};
  \node[dbox, right=of stego, minimum height=13mm] (ch) {Kênh công khai\\(bị giám sát)};
  \node[pbox, right=of ch, minimum height=13mm] (ext) {Bộ trích xuất\\(\tn{extraction})};
  \node[box, right=of ext] (out) {Thông điệp\\khôi phục};
  \draw[arr] (cover.east) -| ($(emb.north west)+(-3mm,0)$) -- ($(emb.north west)+(0,0)$);
  \draw[arr] (msg) -- (emb);
  \draw[arr] (key.east) -| ($(emb.south west)+(-3mm,0)$) -- ($(emb.south west)+(0,0)$);
  \draw[arr] (emb) -- (stego);
  \draw[arr] (stego) -- (ch);
  \draw[arr] (ch) -- (ext);
  \draw[arr] (ext) -- (out);
  \draw[arr, dashed] (key.south) -- ++(0,-3mm) -| node[lbl, pos=0.25, below] {cùng khóa ở phía nhận} (ext.south);
\end{tikzpicture}}
\caption{Mô hình tổng quát của một hệ giấu tin}
\label{fig:stego-model}
\end{figure}

\subsection{Các yêu cầu đối với một hệ giấu tin}

Một hệ giấu tin trong ảnh được đánh giá theo ba yêu cầu có quan hệ đánh đổi lẫn nhau:
\begin{itemize}
  \item \textbf{Tính vô hình} (\tn{imperceptibility}): ảnh mang tin phải giống ảnh phủ đến
        mức mắt người không phân biệt được, thường đo bằng PSNR và SSIM;
  \item \textbf{Dung lượng} (\tn{capacity}): số bit có thể nhúng trên mỗi điểm ảnh; tăng dung
        lượng thường làm giảm tính vô hình;
  \item \textbf{Độ bền vững} (\tn{robustness}): khả năng giữ được thông điệp khi ảnh bị xử lý
        như nén, lọc nhiễu, cắt xén.
\end{itemize}
Ngoài ra còn yêu cầu \textbf{an toàn trước phân tích giấu tin} (\tn{steganalysis}): các phép
kiểm định thống kê không được phân biệt ảnh mang tin với ảnh tự nhiên.

\subsection{Phân loại kỹ thuật giấu tin trong ảnh}

Theo miền nhúng, kỹ thuật giấu tin trong ảnh được chia thành \cite{cheddad2010}:
\begin{itemize}
  \item \textbf{Miền không gian}: sửa trực tiếp giá trị điểm ảnh. Tiêu biểu là kỹ thuật
        \textbf{LSB} (\tn{Least Significant Bit}) thay các bit thấp nhất của điểm ảnh bằng bit
        thông điệp. LSB đơn giản, dung lượng lớn, nhưng rất mong manh trước nén mất dữ liệu.
  \item \textbf{Miền biến đổi}: nhúng vào hệ số của các phép biến đổi như DCT (nền tảng của
        JPEG) hoặc DWT; bền vững hơn nhưng dung lượng thấp hơn.
  \item \textbf{Giấu tin thích nghi}: chọn vị trí nhúng theo độ phức tạp cục bộ của ảnh (vùng
        có nhiều kết cấu) để giảm khả năng bị phát hiện \cite{holub2014}.
\end{itemize}

\subsection{Phân tích giấu tin}

Phân tích giấu tin là lĩnh vực đối nghịch, nhằm phát hiện sự tồn tại của thông điệp ẩn. Với
LSB, hai phương pháp nổi tiếng là \textbf{tấn công thị giác} (quan sát trực tiếp mặt phẳng bit
thấp nhất) và \textbf{kiểm định $\chi^2$ trên cặp giá trị} (\tn{Pairs of Values}) của Westfeld
và Pfitzmann \cite{westfeld2000}. Phép kiểm định thứ hai khai thác hiện tượng: khi LSB bị thay
bằng bit ngẫu nhiên, tần suất của mỗi cặp giá trị $(2k, 2k+1)$ trong histogram bị san bằng.
Phương pháp RS của Fridrich và cộng sự \cite{fridrich2001} còn ước lượng được độ dài thông
điệp đã nhúng. Chương~\ref{chap:vandung} sẽ áp dụng kiểm định $\chi^2$ lên 8000 ảnh mang tin để
đánh giá khả năng bị phát hiện của lược đồ trong bài báo.

\section{Mật mã học và giấu tin: hai lớp bảo vệ bổ sung cho nhau}

Mật mã học và giấu tin giải quyết hai khía cạnh khác nhau của cùng một bài toán.
\bang{tab:crypto-vs-stego} tóm tắt sự khác biệt giữa hai kỹ thuật.

\begin{table}[H]
\centering
\caption{So sánh mật mã học và giấu tin}
\label{tab:crypto-vs-stego}
\small
\begin{tabularx}{\linewidth}{@{}P{3.4cm}LL@{}}
\toprule
\thead{Tiêu chí} & \thead{Mật mã học} & \thead{Giấu tin}\\
\midrule
Mục tiêu chính & Che giấu \emph{nội dung} thông điệp & Che giấu \emph{sự tồn tại} của thông điệp\\
Dạng dữ liệu truyền & Bản mã trông như nhiễu ngẫu nhiên, dễ gây chú ý & Ảnh tự nhiên, không gây nghi ngờ\\
Khi bị phát hiện & Nội dung vẫn an toàn nếu khóa an toàn & Nội dung bị lộ nếu không được mã hóa\\
Cơ sở an toàn & Độ khó tính toán, không gian khóa & Tính vô hình, khả năng chống phân tích thống kê\\
Điểm yếu & Thu hút sự chú ý của kẻ tấn công & Không tự bảo vệ nội dung\\
\bottomrule
\end{tabularx}
\end{table}

Có thể tóm gọn: \emph{mật mã an toàn nhưng không vô hình, còn giấu tin vô hình nhưng không an
toàn} \cite{raut2026hybrid}. Kết hợp hai kỹ thuật theo trình tự ``mã hóa trước, nhúng sau''
cho phép tận dụng ưu điểm của cả hai. Nếu ảnh mang tin không bị nghi ngờ, kẻ tấn công không
biết có gì để tấn công. Nếu bị phát hiện và trích xuất, dữ liệu thu được vẫn chỉ là bản mã
không có ý nghĩa khi thiếu khóa. Thêm vào đó, bản mã có phân bố bit gần đều nên khi nhúng vào
LSB sẽ hành xử như nhiễu ngẫu nhiên, giúp các độ đo chất lượng ảnh có thể dự đoán được bằng
lý thuyết. Đây chính là động lực của khung lai mật mã -- giấu tin được trình bày trong bài
báo và được tái hiện trong báo cáo này.
